
\section{Introduction}
\label{sec:introduction}

\textcolor{red}{
    \textbf{Warning:}
The following paper contains scam content and urls. 
To avoid accidental clicks, we replace all . with [.] of all known malicious urls 
in this paper. 
}

Large language models (LLMs) have become 
critical
infrastructure in software development, with millions 
of developers deploying
AI-generated code for production systems. 
This widespread adoption has occurred
alongside an unprecedented expansion in training data 
scale. Modern LLMs utilize datasets up to 15 trillion tokens,
sourced from web pages, code repositories, and social media
platforms~\citep{openai2023gpt4}. This insatiable demand for
training data has created a fundamental security risk: a large scale
incorporation of malicious content into model weights. 

The internet inherently hosts substantial amounts of 
misinformation, scams, and deliberately poisonous
content~\citep{vosoughi2018spread,lazer2018science,allcott2019trends,
broniatowski2018weaponized, he20244}. 
While traditional web services employ user reporting
mechanisms, and platform-level filtering to combat malicious
material~\citep{gillespie2018custodians,graves2016understanding,
roberts2019behind,roozenbeek2020susceptibility},
the LLM training pipeline operates under a 
fundamentally different paradigm
that amplifies this risk:
data collection prioritizes scale and 
diversity over verification,
crawling billions of pages with minimal quality control.
Once this data is collected, it becomes a training corpus
and is used for training for all future models.
Unlike a search engine that can delist a harmful URL, 
malicious content within a training set is permanently 
embedded into the model's weights. 
Consequently, even if the original source 
 is removed from the web, the malicious data 
 persists and can be  
 replicated across training data of future models, 
 repeatedly exposing end-users to significant harm and risks.


This threat becomes particularly acute in downstream
applications like \textbf{AI-assisted code generation}.
Code generated by LLMs can be integrated into production
systems where it may access sensitive data,
acquire administrative privileges, or cause other direct damage.
Current AI coding assistants can generate thousands of lines
of code in seconds, making it challenging or even impossible for
users to review every line of code generated. 
Moreover, modern software often relies on third-party libraries and
APIs which are very hard for developers to verify every external dependency
used in the generated code. 
A cleverly hidden vulnerability or malicious 
payload can therefore be easily overlooked, 
leading to severe security vulnerabilities unnoticed until
the code is executed, and the damage is done. 
This creates an urgent need to evaluate
the extent to which LLMs are generating malicious
code in practice and to evaluate the potential risks. 
Motivated particularly by a real-world example presented in Section~\ref{sec:example}, 
where a victim lost \$2,500 after ChatGPT generated a code snippet that transmitted his crypto wallet’s private key to a scam URL, in this paper, 
we focus on auditing and evaluating the extent to which
production LLMs generate code containing malicious URLs
in response to completely normal programming prompts that are likely to come from developers. 


% \noindent \textbf{Research Questions:} This paper aims to answer the following
% two key research questions that have profound security implications for the
% software industry in the AI era: 
% \begin{enumerate} 
% \item Are widely deployed production LLMs currently generating
%     malicious code at an alarming, non-negligible rate when facing
%     innocuous programming prompts?

% \item Can we design an automated framework to systematically detect 
%     and expose malicious or poisonous code generated by LLMs
%     at massive scale? 
% \end{enumerate}


\noindent \textbf{Automated Audit Framework:} We develop an automated
audit framework, \framework, to systematically test whether
production LLMs generate code that embeds malicious URLs in
response to developer-style prompts. The key intuition
is that once malicious sources targeting a specific user request exist,
they are rarely isolated; instead, many related variants
also exist which are capable of misleading users toward similarly
harmful outcomes.
When LLMs receive requests for these specific user intents, 
they may reference these malicious variants, 
thereby generating code that contains URLs from scam sites.

Motivated by this observation, our framework begins with
a given seed scam URL and
an oracle capable of detecting malicious URLs. Our framework then
automatically
extracts the context surrounding the 
harmful content in a sandbox,
summarize it, and generate candidate prompts which are normal user 
coding requests.
We subsequently feed these prompts to target production LLMs
for code generation.
Finally, we apply the oracle to identify any generated code snippets
that contain malicious URLs.


This paper focuses on \textit{malicious URLs embedded in code} for two reasons.
First, oracles for malicious URL detection are widely
available and well-established (e.g., Google Safe Browsing~\citep{googleSafeBrowsing},
VirusTotal~\citep{virusTotal}), facilitating large-scale automated
evaluation. Second, malicious URLs in generated code pose severe immediate risks,
ranging from cryptocurrency theft to sensitive data exposure,
making them a high-priority security concern.
Importantly, our automated audit
methodology remains general and can be applied to expose
other forms of
malicious code generation (e.g., backdoors, worms)
provided that appropriate domain-specific
oracles are available.


\noindent \textbf{Results:} \textit{Our experimental results
provide strong empirical evidence that production LLMs can
emit malicious code in response to developer-style prompts
at non-trivial, reproducible rates.}
Through automated auditing of four
production LLMs released in 2024, we find
that on average 4.24\% of code generated in our experiments
contains malicious URLs.
We further constructed a benchmark dataset, \bench, containing 1,377
\textit{developer-style prompts} that trigger
all four LLMs to generate malicious code. \bench is
then applied to seven of the latest production LLMs released in 2025.
We still find that all these models generate malicious code
at a non-negligible rate, ranging from 12.9\% to 47.3\%.
Ranking the models by safety performance yields 
\texttt{Gemini-2.5-Pro} $>$ \texttt{GPT-5} $>$ 
\texttt{Claude-Sonnet-4} $>$ \texttt{Grok-Code-Fast-1} $\approx$ 
\texttt{Gemini-2.5-Flash} $\approx$ 
\texttt{Qwen3-Coder} $\approx$ \texttt{Deepseek-Chat-v3.1}, 
with the top two models exhibiting a statistically significant safety advantage.
These findings demonstrate that
training datasets have been contaminated 
with malicious scam sources at scale.
% , and adversaries, intentionally or not, have successfully poisoned
% training datasets at scale.
Critically, this contamination persists from 2024 models through 2025 releases,
demonstrating that neither current training practices nor safety guardrails
have adequately addressed this vulnerability.
To raise awareness of this urgent threat and support mitigation efforts,
we publicly release our prompts and evaluation results as benchmarks for
future research.


\noindent \textbf{Practical Impact:} While the primary focus of this paper is on the 
developer-style \textbf{prompts} triggering LLMs to generate malicious code that 
contains references to scam phishing sites, 
we have unexpectedly discovered numerous scam phishing sites from our auditing framework 
that were not included in the scam databases we used, 
yet are generated by LLMs. 
This is particularly impressive, as the knowledge cutoffs of 
these LLMs all predate August 2024 as shown in Table~\ref{tab:llm_specifications}, 
indicating that these malicious sites have likely been active for 
over a year while evading detection by conventional security measures. 
We have reported all of these sites to the respective scam database maintainers. 
As of this paper's submission, 
\textbf{62} of these sites have been confirmed and 
added to the major scam database 
eth-phishing-detect~\cite{MetaMask_eth-phishing-detect_2025}, with 
proofs available at \url{https://sites.google.com/view/scam2prompt}.


\noindent \textbf{Contributions:} This paper makes the following contributions:

\begin{itemize} [nosep,leftmargin=*]
        \item \textbf{Empirical Evidence of Malicious Code generated by
            Production LLMs:} We disclose and evaluate the extent to which
            production LLMs can generate malicious code, demonstrating that at
            4.24\% of LLM-generated code contains malicious URLs alone when
            responding to our developer-style prompts. The actual rate of malicious
            code generation likely exceeds this figure when considering attack
            vectors beyond URLs.

        \item \textbf{\framework: Automated Auditing Framework:} We
            design and implement a scalable framework that given seed
            malicious sources and domain-specific oracles, automatically
            generates prompts appearing as developer-style coding requests
            while systematically exposing malicious code generation in
            production LLMs. This general methodology applies to any type of
            malicious behavior provided appropriate oracles are available.
            To support reproducibility, we release the source code of \framework at
            \url{https://github.com/anonymous00002/Scam2Prompt-ICML-Artifact/}.

        \item \textbf{\bench:} We release a selected benchmark of 1,377
            developer-style prompts that trigger all four 
            production LLMs audited by \framework 
            to generate malicious code. We further apply
            this benchmark to seven latest production LLMs released
            in 2025, demonstrating that all models generate malicious code
            at a non-negligible rate ranging from 12.9\% to 47.3\%.
            To support reproducibility, we release the dataset at
            \url{https://huggingface.co/datasets/anonymous-author-32423/Innoc2Scam-bench/tree/main}, 
            and a tutorial at \url{https://huggingface.co/datasets/anonymous-author-32423/How-To-Use-Innoc2Scam-Bench}. 

        \item \textbf{Defense Evaluation:} We evaluate \bench
            against two widely used defenses (guardrails and RAG-based agents)
            and show they provide limited risk reduction, 
            further demonstrating the need for adding  
            explicit oracle-based URL checks after code generation.

            
\end{itemize}
